\subsection{Reproducible Settings and Timing}
The experiments use the fixed MovingAI \texttt{random-32-32-10} map, with 32 by 32 cells and one metre per cell. Robot actions are the four cardinal moves or waiting; each control cycle is one second. A square FoV of radius seven cells and five future intention cells are retained. The UAV starts at the map centre and its analytical coverage-optimal reference altitude. Its three-dimensional motion is limited to one metre per control cycle toward the maximum-coverage target computed from successfully received position reports, with a 20 m minimum altitude. No fresh simulator position is supplied to the allocation or placement controllers.

The clean frozen PANAMA checkpoint qualified on 200 episodes at each population: success rates were 97\%, 96\%, 95.5\%, and 95\% for 2, 4, 6, and 8 robots. Evaluation uses 50 separate held-out cases per population and a 250-step cap. All treatments share policy weights, starts, goals, seeds and orbital offsets. Hashes identify the checkpoint, map, case bank, channel archive and source code. Map compatibility permits exact LF/CRLF conversion only.

The 100 ms control observation budget determines which cached information is available before an action. It is not a hard task deadline. One logical coordination task is generated per robot per cycle after the initial fresh briefing. The task size and CPU demand below are simulation assumptions, not profiled neural-inference workloads or measured processor performance. Whole tasks retain their execution site once assigned; the finite processors account for existing work and recompute priority-aware sharing on arrivals and completions. CPU service is work-conserving in the latency-only case and may intentionally idle under the DVFS objective. Actual dynamic CPU energy integrates $\kappa_D f_i(t)^3$ over service.

\begin{table}[t]
\centering\caption{Nominal access, computation and timing assumptions.}
\label{tab:nominal_parameters}\scriptsize
\begin{tabular}{@{}p{0.58\columnwidth}p{0.35\columnwidth}@{}}
\toprule Parameter & Value\\\midrule
G2A carrier; per-link bandwidth & 2.5 GHz; 20 MHz\\
Urban $a,b$; excess LoS/NLoS losses & 9.61, 0.16; 1/20 dB\\
Thermal PSD; receiver noise figure & $-174$ dBm/Hz; 7 dB\\
Effective receiver PSD & $-167$ dBm/Hz\\
Robot conducted uplink power & 0.2 W\\
Access uplink/downlink service ceilings & 50/50 Mbit/s\\
Task logical input; cycles per bit & $10^5$ bits; $10^3$\\
Drone/cloud CPU capacity & $10^9/10^{10}$ cycles/s\\
Drone dynamic CPU coefficient & $\kappa_D=10^{-28}$\\
Resource weights; completed-agent weight & $[1,4]$; 1\\
Latency/energy coefficients & $1$ s$^{-1}$/$1$ J$^{-1}$\\
Control cycle; observation budget & 1 s; 0.1 s\\
Maximum fragment; wire overhead & 65,536 B; 48 B/fragment\\
Hard mission-task deadline & None\\\bottomrule
\end{tabular}
\end{table}

\subsection{Channels, LEO Geometry and Propagation}
G2A uses the probabilistic urban LoS mean path loss in \eqref{eq:g2a_pathloss}; no instantaneous small-scale fading is drawn. The thermal PSD and receiver noise figure are applied once, so the effective PSD equals $-167$ dBm/Hz. The rate in \eqref{eq:g2a_rate} is additionally capped by the configured shared access-service capacity. Access is time-shared with round-robin queue service; the bandwidth is not multiplied by the number of simultaneous robots. Upload input bits and actual serialized wire bits are recorded separately. To retain checkpoint compatibility, the policy's third observation channel remains PANAMA's normalized coverage-margin proxy. That proxy is distinct from the physical RF SNR used to calculate the actual access rate and is not reported as a calibrated SNR measurement.

The orbital archive contains six hours of one-second SGP4/Skyfield geometry replayed at its archived epoch, 31 August 2026. The reference site is $(46.2^{\circ},-74.6^{\circ})$ at 440 m elevation; the assumed gateway is $(45.3^{\circ},-74.6^{\circ})$ at 300 m. These coordinates are model inputs, not a measured provider gateway. The highest visible satellite is selected at 15 s scheduling boundaries, with a $25^{\circ}$ terminal and $10^{\circ}$ gateway elevation mask. Original 50 ms attachment/reconfiguration interruptions are retained. Frames are 20 ms with a 5\% service guard.

To implement the association-average assumption, terminal and gateway ranges and statistical attenuation in dB are averaged within each association/scheduling interval. Free-space loss, SNR and Shannon rates are then recomputed from those mean inputs; arithmetic means of rates are not substituted. ITU-Rpy 0.4.0 provides the archived P.618 statistical atmospheric predictions \cite{iturP618} at a fixed 1\% exceedance quantile. This quantile is not a sampled weather process. No provider modulation or beam-scheduler reconstruction is claimed. Small within-map UAV displacement is neglected in the orbital geometry; the conservative round-trip propagation error bound is below $0.7$ microseconds, while G2A uses the actual UAV placement.

\begin{table}[t]
\centering\caption{Backhaul engineering assumptions and full-route delays.}
\label{tab:backhaul_parameters}\scriptsize
\begin{tabular}{@{}p{0.58\columnwidth}p{0.35\columnwidth}@{}}
\toprule Parameter & Value\\\midrule
Uplink/downlink carrier & 14.25/11.7 GHz\\
Uplink/downlink bandwidth & 20/100 MHz\\
Drone conducted satellite-uplink power & 2 W\\
Terminal EIRP; implied uplink antenna gain & 35 dBW; 31.99 dBi\\
Satellite receive $G/T$ & $-5$ dB/K\\
Satellite EIRP; terminal receive $G/T$ & 30 dBW; 5 dB/K\\
Implementation loss; Shannon efficiency & 3 dB; 0.5\\
Atmospheric antenna diameter; exceedance & 0.3 m; 1\%\\
Nominal backhaul uplink/downlink ceilings & 10/50 Mbit/s\\
One-way ground-network propagation & 10 ms\\
Association-average full-route RTT & 24.92--29.60 ms\\\bottomrule
\end{tabular}
\end{table}

The full route is robot--UAV--LEO--gateway--ground cloud and its reverse, with no inter-satellite hop. Each direction is limited by both its terminal and gateway link budgets. Conducted transmitter power, rather than EIRP, determines transmission energy. Propagation includes one outbound and one return contribution; fragmentation does not multiply propagation by the number of fragments. Actual response serialization is simulated even though it is omitted from the analytical latency expression. Solver wall time is measured separately from logical task CPU time.

\subsection{Validation, Baselines and Metrics}
Six predictable sanity cases cover large satellite delay, fast cloud/slow drone, unavailable backhaul, unavailable drone CPU, fixed-route priority differentiation and a DVFS optimum with idle CPU. Every expected assignment is checked using an independent primal CVXPY formulation and exhaustive enumeration. The earlier supplied four-robot MATLAB validation inputs are also preserved, including the original infinite deadlines and the separate Case E feasibility test; execution here uses the approved CVXPY equivalent.

The primary grid combines uplink and common CPU multipliers $\{0.25,1,4\}$ under equal resource weights, episode-initial static priorities, AoI-only weights and current PANAMA-priority weights, with one fresh-coordination reference. Separate sensitivity conditions vary drone CPU, cloud CPU, downlink, logical input and cycles/bit by 0.25 and 4, the weight ratio by 1 and 16 (baseline 4), and the energy coefficient by 0, 0.1 and 10 (latency coefficient 1). There are 97 total conditions and 19,400 paired episodes. The ideal fresh reference has instantaneous coordination and no modeled support-task energy; its energy is not a hardware benchmark.

Task outcomes include observed and snapshot-predicted latency, prediction error, the 95th latency percentile, unfinished-task counts, stage service times, and modeled energy by transmitter and processor. Mission outcomes include completion, capped mission duration, successful makespan, sum of time-to-first-arrival costs, stationary steps (including blocked moves), and collision attempts. Unfinished missions remain censored at 250 steps; PAR-2 assigns them 500. Directed inconsistency is recorded with UL/DL/E2E decomposition, square-FoV relevance, source-time AoI and raw numerator/denominator counts. Mean per-episode inconsistency and pooled relevant-pair ratios are labelled separately. A zero relevance denominator gives zero.

Each comparison uses the same held-out case in both conditions. We report paired mean differences and seeded 2,000-resample 95\% bootstrap intervals, with the number of usable pairs for metrics affected by censoring. The intervals are pointwise across the sensitivity grid and are not family-wise multiple-comparison guarantees. Conditional successful makespans are secondary to completion and PAR-2. Snapshot optimality and a lower relaxed objective do not imply an improvement in every mission outcome.

\subsection{Measured Comparisons}
The complete paired study is in progress. Numerical result tables and figures are generated from the finished episode records; no placeholder values are substituted.

